\section{Spiking Neural Network (SNN) Architecture}
\label{sec:snn_architecture}

The core focus, effort and innovation of this project is the development and optimization of the \texttt{SNN\_MAX\_REPLICA\_SE} architecture. 
This model is explicitly designed to process high-definition neuromorphic event data while adhering to the computational paradigms of \gls{snn}s. 
By prioritizing event-driven, binary communication, the architecture aims to drastically reduce both the dynamic power consumption and the computational complexity compared to traditional deep learning models, making it superior for power and time constrained environments.

\subsection{Spatiotemporal Input \& Multi-Step Execution}
\label{subsec:snn_input}
The raw \gls{etram} event data is initially preprocessed to a spatial resolution of $640 \times 360$ with $T=10$ discrete temporal bins. 
Consequently, the network ingests a massive, high-dimensional spatiotemporal tensor of shape $[B, 10, 2, 360, 640]$ (Batch, Time, Polarity, Height, Width). 

Given the high spatial dimensionality and the iterative nature of spiking neurons, native Python loop execution over the time dimension proved computationally restrictive. 
To overcome this bottleneck, the entire architecture was refactored to utilize the SpikingJelly framework's optimized Multi-Step (\texttt{step\_mode='m'}) execution paradigm \cite{fang2023spikingjelly};
instead of iterating through the 10 time bins sequentially, the entire spatiotemporal tensor is passed directly into custom-compiled CUDA kernels.
This low-level optimization enables high-speed, parallelized execution of the temporal dynamics, making the extensive benchmarking of this HD dataset feasible on local GPU hardware, and reducing overall training time from approximately 16 hours to just 7.

\subsection{Neuron Dynamics: \gls{lif}}
\label{subsec:snn_lif}
At the foundation of the architecture lies the \gls{lif} neuron model as defined in Section \ref{subsec:bg_lif}, which governs the temporal dynamics and binary spike generation of the network. 

\subsubsection{The Forward Pass (Spike Generation)}
Both the membrane potential dynamics from Equation \ref{eq:lif_membrane} and Heaviside spike function from Equation \ref{eq:lif_spike} are configured with the following project-specific parameters:

\begin{align*}
    V_{th} &= 0.2 \\
    \tau &= \langle 4.0 \rangle
\end{align*}

The voltage threshold $V_{th}$ was lowered to 0.2 to enable more information flow through the network; in earlier revisions, this was set to 0.3, which caused the network to become silent, whereas decreasing it further would result in too much noise.
The time constant $\tau = 4.0$ creates a leaky effect, where the neuron remembers information for roughly four to five timesteps.
In deep spiking networks, a lower $\tau$ value can cause the signal to die before it reaches the regression head, since the neuron forgets the input too quickly; in this implementation, $\tau = 4.0$ was determined to be the optimal value where the input was not forgotten, whilst preventing over-saturation.

\subsubsection{The Backward Pass (Surrogate Gradients)}

As described in Section \ref{sec:background_ann}, the Heaviside function's non-differentiable property poses a challenge in spiking neural networks.
To get around the zeroed gradients and the consequent ``dead neuron'' problem, the \gls{snn} model employs the \gls{atan} function described in \ref{subsec:bg_surrogate},
where the backward pass approximates the step function using Equation \ref{eq:atan_function};
this mathematical substitution allows the network to learn complex spatiotemporal representations from scratch while maintaining a purely binary forward execution.
In this implementation, SpikingJelly's default value of $\alpha=2.0$ was used, which is the standard value used for high-performance \gls{snn} training.

\subsection{Backbone: SEW-ResNet \& Attention}
\label{subsec:snn_backbone}

To extract high-level semantic features and manage the computational complexity of the high-definition input, the architecture utilizes a deep residual backbone divided into five sequential stages, as outlined in Table \ref{tab:snn_staged_architecture}.

\begin{table}[htbp]
\centering
\resizebox{\textwidth}{!}{%
\begin{tabular}{|c|l|c|c|p{6cm}|}
\hline
\textbf{Stage} & \textbf{Name} & \textbf{Resolution} & \textbf{Channels} & \textbf{Primary Operation} \\
\hline
1 & Initial Stem      & 90$\times$160 & 64  & Stride-2 Conv + \gls{lif} (Downsamples HD input) \\
2 & Backbone Alpha    & 22$\times$40  & 256 & Blocks 1 \& 2 (SEWResBlocks with Stride-2) \\
3 & Attention Bridge  & 22$\times$40  & 256 & Squeeze-and-Excitation (SE) Global Firing Recalibration \\
4 & High-Res Backbone & 22$\times$40  & 512 & Blocks 3 \& 4 (SEWResBlocks with Stride-1) \\
5 & Temporal Fusion   & 22$\times$40  & 512 & Summation of spikes across all $T$=10 timesteps \\
\hline
\end{tabular}%
}
\caption[Network Architecture Stages]{Detailed breakdown of the five sequential stages in the \texttt{SNN\_MAX\_REPLICA\_SE} backbone.}
\label{tab:snn_staged_architecture}
\end{table}

\subsubsection{Spatial Downsampling}
As illustrated, the architecture implements an aggressive $16\times$ total spatial downsampling strategy to compress the initial $640 \times 360$ input. 
This is achieved through strided convolutions applied within the Initial Stem and Backbone Alpha, progressively reducing the spatial dimensions to a highly compressed $22 \times 40$ grid while expanding the channel width up to 512.

\subsubsection{Spike-Element-Wise (SEW) Residual Blocks}
Building deep \glspl{snn} presents significant challenges due to the ``vanishing spike'' problem inherent in standard residual connections. 
In a standard \gls{resnet}, the residual addition occurs \textit{before} the activation function.
However, if we follow this standard and add two binary spike trains, the result is a non-binary, continuous value (e.g., $1+1=2$), which destroys the energy-efficient, event-driven nature of the network \cite{fang2022resnet_sew}.

As detailed in Section \ref{subsec:bg_sew}, and visually in Figure \ref{fig:sewresblock_bg}, \gls{sew} ResBlocks solve
the vanishing spike problem by placing the residual operation after the spiking activation function \cite{fang2022resnet_sew}.
This architecture exploits this property by stacking 4 \texttt{SEWResBlock} units across Stages 2 and 4 of the backbone (Table \ref{tab:snn_staged_architecture}),
enabling representational depth that would otherwise induce spike degradation under a standard residual equivalent.

\subsubsection{Structural Attention: SE-Block}
\label{subsec:snn_se_block}

In deep learning, a bottleneck refers to a point in the neural network where the information flow is forced through its most compressed and abstract representation. 
In this \texttt{SNN\_ULTRA} architecture, the bottleneck exists at the transition point between stages 2 and 3 of the backbone, and represents two things:

  \begin{enumerate}
      \item \textbf{A Spatial Constraint (Resolution Limit):} 
      At this stage, the original $180 \times 320$ spatiotemporal input has been downsampled by a factor of 8x via the strided convolutions in the stem and early residual block, resulting in a highly compressed$22 \times 40$ feature map.
      This spatial bottleneck is critical because each remaining "pixel" must now represent a large $8 \times 8$ receptive field of the original HD sensor data, forcing the network to discard fine-grained geometric redundancy while retaining the
      core structural boundaries of multiple detected objects.

      \item \textbf{A Semantic Shift (High-Level Abstraction):} 
      This juncture marks the definitive pivot where the network transitions from extracting low-level geometric primitives, such as spike-edge orientation and motion direction, to synthesizing complex semantic concepts. 
      Without an explicit selection mechanism, the network at this depth is susceptible to noise and artifacts from the HD event stream, which can blur the distinction between legitimate targets and background interference.
  \end{enumerate}

To both mitigate and solve these bottleneck constraints, the \gls{se} mechanism depicted in Section \ref{subsec:bg_se} and Figure \ref{fig:se_block_bg} is introduced in Stage 3 of the backbone to operate on the $22\times40$ feature map, sitting directly at the transition point \cite{hu2019_se_block}.
The combination of the ``Squeeze'' and ``Excitation'' phases effectively mitigate the spatial constraint by allowing the network to incorporate global context into its local representations, compensating for the loss of resolution.

When applied in the context of an \gls{snn}, this acts as a dynamic firing modulator: the backbone's activations are re-calibrated by looking at the global context of the entire frame, magnifying the most informative channels whilst suppressing those containing artifacts.
By applying this recalibration, the \gls{se} block ensures that only the most relevant ``Objectness`` (the probability that a specific region of an image contains a generic object of interest rather than background noise) features are passed into the final high-capacity blocks, effectively mitigating the semantic ambiguity before the final bounding box regression.

\subsection{Firing Rate Regularization (FRR)}
\label{subsec:snn_frr}

As detailed in Section \ref{subsec:bg_frr}, an \gls{frr} metabolic penalty (Equation \ref{eq:frr}) is integrated into the training loss to prevent
network silencing and over-saturation \cite{roy2019towards}.
For the \texttt{SNN\_ULTRA} architecture, the penalty is configured with $\lambda = 5.0$ and a target firing rate of $\text{Target} = 20\%$.
Using a high $\lambda$ value punishes the model for being silent, forcing the backbone to wake up and fire.
Setting the target firing rate to $20\%$ was chosen designed to balance
spatiotemporal integration with power efficiency - providing enough temporal resolution to
draw smooth bounding boxes while still maintaining $80\%$ sparsity for energy savings.

\subsection{Multi-Box Regression Head}
\label{subsec:snn_head}
Following the final backbone stage, the abstract spatiotemporal feature maps are fed into a unified detection head. 
To process the spatial information effectively and mitigate spatial ambiguities observed during the iterative design process (detailed in Section \ref{sec:architectural_evolution}), the head utilizes an $8 \times 8$ Adaptive Average Pooling mechanism. 

This pooling operation aggregates the features into 64 distinct spatial receptive blocks. 
The pooled features are then mapped via final convolutional and linear layers to an output tensor of shape $[B, 5, 5]$, where $B$ represents the batch size.
The final two dimensions encode the top 5 predicted bounding boxes globally per image, with each box defined by a feature tuple $\mathcal{F} = \{x, y, w, h, c\}$. 
Here, $(x, y, w, h)$ represent the normalized spatial coordinates and dimensions of the bounding box, and $c$ denotes the object confidence logit. 
This formulation allows the network to predict multiple vehicles simultaneously, which are subsequently evaluated using a vectorized Multi-Box Loss function.

\todo[inline,backgroundcolor=myblue,textcolor=white]{Add model architecture diagram.}
\todo[inline,backgroundcolor=myblue,textcolor=white]{Check \& Add model hyperparameter table.}